\chapter{Iron Needs}
\index{Iron Needs}

\section*{Definition and Overview}
Iron is an essential mineral that the body needs to make haemoglobin, the protein in red blood cells that carries oxygen, as well as myoglobin in muscle and enzymes used throughout the body. How much iron a person needs depends on age, sex, growth, pregnancy, and blood loss, and needs change across the lifespan. This chapter explains recommended daily iron intakes for different groups, who is at risk of deficiency, and how to meet iron needs through food and, where necessary, supplements. Understanding your own iron requirements is the first step in preventing deficiency and the fatigue and other symptoms it causes.

\section*{Epidemiology}
Iron deficiency is the most common nutritional deficiency in the world and affects around one in eight women of reproductive age. Menstruating women need almost twice as much iron as men because of monthly blood loss, and pregnant women have the highest needs of all, rising to around 27 milligrams daily. Babies and toddlers need relatively high amounts for growth, and teenagers, especially girls, are often deficient. Vegetarians and vegans are at risk because plant iron is absorbed less efficiently than the iron in meat. Older people and those with gut conditions or chronic disease also have increased needs or impaired absorption.

\section*{Aetiology and Pathophysiology}
Iron needs are determined by three factors: the iron required for growth and red blood cell production, the iron lost from the body, and the efficiency of iron absorption. Most iron in the body is recycled, but menstruation, pregnancy, growth, and bleeding create real losses that must be replaced from the diet. The body absorbs haem iron from meat, poultry, and fish more efficiently than non-haem iron from plants, and absorption is increased by vitamin C and decreased by tea, coffee, calcium, and phytates in grains and legumes. When needs outstrip intake, the body first depletes its iron stores and then limits red blood cell production, producing iron deficiency and anaemia.

\section*{Clinical Features}
\begin{itemize}
    \item Recommended daily intakes vary by group: around 8 milligrams for adult men, 18 milligrams for menstruating women, and 27 milligrams in pregnancy
    \item Babies from six months, toddlers, and teenagers have high relative needs
    \item Women who breastfeed need about 9 to 10 milligrams daily
    \item People with heavy periods, frequent blood donation, or gut bleeding need more
    \item Vegetarians and vegans may need up to 80 per cent more iron than meat-eaters
    \item Adequate iron prevents the fatigue, paleness, and breathlessness of deficiency
    \item Iron needs are best met by a varied diet, with supplements used when needs cannot be met by food
\end{itemize}

\section*{Differential Diagnosis}
When iron status is uncertain, the cause of any deficiency must be identified.
\begin{itemize}
    \item \textbf{Dietary inadequacy:} low iron intake, especially in vegetarians and vegans
    \item \textbf{Blood loss:} heavy periods, gut bleeding, and repeated blood donation
    \item \textbf{Malabsorption:} coeliac disease, inflammatory bowel disease, and gastric surgery
    \item \textbf{Increased demand:} pregnancy, growth, and chronic illness
    \item \textbf{Low absorption:} high tea, coffee, or phytate intake with meals
    \item \textbf{Other causes of anaemia:} B12 or folate deficiency, thalassaemia, and chronic disease
\end{itemize}

\section*{When to Seek Medical Care}
Anyone with fatigue, paleness, or other symptoms of anaemia should have their iron checked, and people with heavy periods, gut symptoms, or restricted diets should discuss their iron needs with a doctor. Pregnant women are screened and advised about supplementation. Taking iron supplements without a confirmed need can cause harm, so advice should be sought before starting.

\textbf{Contact your local emergency services for:}
\begin{itemize}
    \item Severe breathlessness, chest pain, or collapse
    \item Black, tarry, or bloody stools, or vomiting blood
    \item Heavy bleeding with faintness or dizziness
    \item Sudden severe weakness or confusion
\end{itemize}

\section*{Diagnosis and Investigations}
\begin{itemize}
    \item \textbf{Full blood count:} haemoglobin and red cell indices
    \item \textbf{Ferritin:} the measure of iron stores, which falls before anaemia develops
    \item \textbf{Iron studies:} serum iron and transferrin saturation where needed
    \item \textbf{Dietary assessment:} review of iron intake against needs
    \item \textbf{Assessment of blood loss:} menstrual history, stool tests, and gut investigation where indicated
    \item \textbf{Screening in pregnancy:} routine iron checks in antenatal care
    \item \textbf{Follow-up testing:} confirming that stores are replenished after treatment
\end{itemize}

\section*{Management}
\begin{itemize}
    \item \textbf{Dietary intake:} meeting needs with iron-rich foods, combining plant iron with vitamin C
    \item \textbf{Iron supplements:} for people who cannot meet needs by diet, such as pregnant women
    \item \textbf{Treatment of deficiency:} oral iron, and intravenous iron for intolerance or malabsorption
    \item \textbf{Treatment of the cause:} managing heavy periods, gut bleeding, or absorption problems
    \item \textbf{Education:} advice on cooking, food combinations, and avoiding inhibitors
    \item \textbf{Monitoring:} repeat blood tests to confirm correction
    \item \textbf{Public health measures:} iron-fortified foods and infant feeding advice
\end{itemize}

\section*{Complications}
\begin{itemize}
    \item Iron deficiency and iron deficiency anaemia
    \item Fatigue, impaired concentration, and reduced work capacity
    \item Impaired cognitive development in infants and children
    \item Adverse pregnancy outcomes, including preterm birth and low birth weight
    \item Restless legs syndrome
    \item Iron overload from unnecessary supplementation
    \item Missed underlying disease when blood loss is not investigated
\end{itemize}

\section*{Prognosis and Outcomes}
When iron needs are understood and met, deficiency is preventable and correctable. Symptoms respond within weeks of treatment, and blood counts normalise within one to two months. Meeting needs through a varied diet is achievable for most people, and targeted supplementation protects those at highest risk, including pregnant women and infants. The key is recognising that needs vary across the lifespan and adjusting intake accordingly.

\section*{Prevention and Self-Care}
\begin{itemize}
    \item Know your iron needs and whether you are in a high-risk group
    \item Eat iron-rich foods: red meat, poultry, fish, legumes, eggs, and fortified cereals
    \item Combine plant iron with vitamin C from fruit and vegetables
    \item Avoid tea and coffee with iron-rich meals
    \item Take iron supplements in pregnancy as advised
    \item Have blood tests if you have persistent fatigue
    \item Feed babies iron-rich or iron-fortified foods from six months
    \item See a doctor before starting iron supplements without a known deficiency
    \item Treat heavy periods and seek medical advice for ongoing blood loss
\end{itemize}

\section*{Sources and Further Reading}
The following reputable sources provide further information for healthcare students and patients seeking reliable, current advice about iron needs.
\begin{itemize}
    \item National Health and Medical Research Council. \textit{Australian dietary guidelines and nutrient reference values.} https://www.nhmrc.gov.au/
    \item Healthdirect Australia. \textit{Iron and iron deficiency.} https://www.healthdirect.gov.au/
    \item Dietitians Australia. \textit{Iron in the diet.} https://dietitiansaustralia.org.au/
    \item Pregnancy, Birth and Baby. \textit{Iron in pregnancy.} https://www.pregnancybirthbaby.org.au/
    \item NHS. \textit{Vitamins and minerals: iron.} https://www.nhs.uk/
    \item World Health Organization. \textit{Iron deficiency and anaemia.} https://www.who.int/
\end{itemize}
